\chapter{Ejercicios}

\section{Sección 1}
% //TODO: ejercicios hacer 1.1 - 1.6 ? - 1.11 - 1.13 - 1.15? 


\section{Sección 2}
% //TODO: ejercicios hacer 2.2 - 2.5* - 2.6 - 2.8 - 2.12* - 2.16?


\begin{ejercicio}[Libro 2.10]
    % //TODO: copiar ejercicio de mis apuntes y Cristian
\end{ejercicio}

\section{Sección 3}
% //TODO: ejercicios hacer 3.3* - 3.6 - 3.8 -  3.10 - 3.11 - 3.14 - 3.17 - 3.18 - 3.19 - 3.20 b HACER - 3.21*


\begin{ejercicio}[Libro 3.6]
    Sea $(\Omega, \mathcal{F}, P)$ un espacio de probabilidad y $A, B, C \in \mathcal{F}$. 
    Demuestre directamente (sin usar el Corolario (3.20) ni el Teorema (3.24)):
    
    \begin{enumerate}
        \item[(a)] Si $A$ y $B$ son independientes, entonces también lo son $A$ y $B^c$.
        \item[(b)] Si $A, B, C$ son independientes, entonces también lo son $A \cup B$ y $C$.
    \end{enumerate}
    Haremos el ejercicio por partes, veámoslo.
    \begin{itemize}
        \item [(a)] Sabemos que $A$ y $B$ son independientes, por lo que se cumple que $P(A \cap B) = P(A) \cdot P(B)$, veámos que tenemos
        \begin{equation*}
            P(A \cap B^c) = P(A) - P(A \cap B) = P(A) - P(A) \cdot P(B) = P(A) \cdot (1 - P(B)) = P(A) \cdot P(B^c)
        \end{equation*}
        donde hemos usado la independencia de $A$ y $B$.
        \item [(b)] Sabemos que $A, B, C$ son independientes, por lo que se cumple que $P(A \cap B \cap C) = P(A) P(B) P(C)$, veámos que tenemos
        \begin{align*}
            & P((A \cup B) \cap C) = P((A\setminus B \cup B) \cap C)= P((A \setminus B \cap C) \cup (B \cap C)) = P(A \setminus B \cap C) + P(B \cap C) = \\
            & = P(A \cap B^c \cap C) + P(B \cap C) = P(A)P(B^c)P(C) + P(B)P(C) = \\
            & = P(A)(1 - P(B))P(C) + P(B)P(C) = P(C)[P(A)P(B^c) + P(B)] \overset{(*)}{=} P(C)P(A \cup B)
        \end{align*}
        donde en $(*)$ hemos usado lo siguiente
        \begin{align*}
            & P(A)P(B^c) =P(A\cap B^c)=P(A\setminus B) \\
            & P(A)P(B^c) + P(B) = P(A\setminus B) + P(B) = P((A\setminus B) \cup B) = P(A \cup B)
        \end{align*}
    \end{itemize}
\end{ejercicio}


\begin{ejercicio}[Libro 3.10]
    % //TODO: copiar ejercicio de mis apuntes, Jesus y Cristian
\end{ejercicio}

\begin{ejercicio}[Libro 3.17 Thinning of a Poisson distribution]
    Supongamos que el número de huevos $N$ que pone un insecto sigue una 
    distribución de Poisson con parámetro $\lambda > 0$:
    $$N \sim \text{Poisson}(\lambda)$$
    De cada huevo nace una larva con probabilidad $p \in (0,1)$, de forma 
    independiente para cada huevo. Encuentre la distribución del número 
    total de larvas $j$.

    Para la resolución definiremos la v.a. $N$ que representa el número de huevos que pone el insecto, de manera que tenemos
    \begin{equation*}
        N: \Omega \to \mathbb{N}_0, \quad N(w)= n \text{ si el insecto pone } n \text{ huevos en el experimento } w \in \Omega
    \end{equation*}
    es importante destacar que el espacio muestral $\Omega$ es el conjunto de todos los posibles experimentos que podemos realizar, y la variable aleatoria $N$ asigna a cada experimento el número de huevos que pone el insecto en dicho experimento. Además, sabemos que $N$ sigue una distribución de Poisson con parámetro $\lambda > 0$, es decir, tenemos que
    \begin{equation*}
        P(N=k) = e^{-\lambda} \frac{\lambda^k}{k!}, \quad k \in \mathbb{N}_0
    \end{equation*}
    Por otra parte, definimos los eventos $A_i$ como el evento en el que nace una larva del huevo $i$-ésimo, de modo que tenemos la variable aleatoria $S$ que representa el número de larvas que nacen en un experimento $w\in \Omega$, es decir,
    \begin{equation*}
        S: \Omega \to \mathbb{N}_0, \quad S(w) = \sum\limits_{i=1}^{N(w)}\mathbbm{1}_{A_i}(w) \ \forall w \in \Omega
    \end{equation*}
    donde $\mathbbm{1}_{A_i}$ es la función indicadora del evento $A_i$, es decir,
    \begin{equation*}
        \mathbbm{1}_{A_i}(w) = \begin{cases}
            1, & \text{si nace una larva del huevo } i \text{-ésimo en el experimento } w \in \Omega \\
            0, & \text{si no nace una larva del huevo } i \text{-ésimo en el experimento } w \in \Omega
        \end{cases}
    \end{equation*}
    donde tenemos que $P(A_i) = p$ para todo $i \in \mathbb{N}$, y los eventos $A_i$ son independientes entre sí. \\ \\
    \noindent
    Nuestro objetivo es encontrar la distribución de la variable aleatoria $S$, es decir, encontrar $P(S=j)$ para todo $j \in \mathbb{N}_0$. Para ello primero \underline{fijaremos} $n\in \N_0$ como el número de huevos que pone el insecto en un experimento $w \in \Omega$, es decir, $N(w) = n$, y definiremos la siguiente medida de probabilidad para $S_n$, que representa el número de larvas que nacen cuando el insecto pone $n$ huevos, que sigue claramente una distribución binomial con parámetros $n$ y $p$
    \begin{equation*}
        P(S_n=j) = \binom{n}{j} p^j (1-p)^{n-j}, \quad \forall j\leq n, \ j \in \mathbb{N}_0
    \end{equation*}
    Ahora sí, aplicando la fórmula de la probabilidad total, tenemos que
    \begin{align*}
        & P(S=j)  = \sum\limits_{k=j}^{\infty} P(S=j | N=k) P(N=k)\overset{(*)}{=} \sum\limits_{k=j}^{\infty} P(S_k=j) P(N=k) = \sum\limits_{k=j}^{\infty} \binom{k}{j} p^j (1-p)^{k-j} e^{-\lambda} \frac{\lambda^k}{k!} = \\
        & = e^{-\lambda} \frac{p^j}{j!} \sum\limits_{k=j}^{\infty} \frac{(1-p)^{k-j} \lambda^k}{(k-j)!} = e^{-\lambda} \frac{p^j}{j!} \lambda^j \sum\limits_{m=0}^{\infty} \frac{((1-p)\lambda)^m}{m!} \overset{(**)}{=} e^{-\lambda} \frac{(p\lambda)^j}{j!} e^{(1-p)\lambda} = e^{-p\lambda} \frac{(p\lambda)^j}{j!}
    \end{align*}
    donde en $(**)$ hemos usado la serie de Taylor de la función exponencial, y en $(*)$ hemos de usado el siguiente desarrollo
    \begin{align*}
        P(S=j\mid N=k)=P(\sum\limits_{i=1}^{k} \mathbbm{1}_{A_i} = j\mid N=k)=\dfrac{P(\{\sum\limits_{i=1}^{k} \mathbbm{1}_{A_i} = j\}\cap \{N=k\})}{P(N=k)}\overset{(***)}{=}P(\sum\limits_{i=1}^{k} \mathbbm{1}_{A_i} = j)=P(S_k=j)
    \end{align*}
    donde se ha usado en $(***)$ que el número de huevos $N$ es independiente del nacimiento de las larvas en cada huevo, es decir, $A_1,\ldots, A_k$. Finalmente se tiene que la distribución del número total de larvas $j$ es una distribución de Poisson con parámetro $p\lambda$, es decir,
    \begin{equation*}
        S \sim \text{Poisson}(p\lambda) \quad \text{ donde } P(S=j) = e^{-p\lambda} \frac{(p\lambda)^j}{j!}, \quad j \in \mathbb{N}_0
    \end{equation*}
\end{ejercicio}


\begin{ejercicio}[Libro 3.20]
    % //TODO: copiar ejercicio de Jesus y Cristian
\end{ejercicio}

\section{Sección 4}
% //TODO: ejercicios hacer 4.15 b)

\begin{ejercicio}[Libro 4.15 b)]
    % //TODO: copiar ejercicio de Jesus y Cristian
\end{ejercicio}

\section{Sección 5}
% //TODO: ejercicios hacer 5.1 - 5.2 - 5.4 - 5.10* - 5.17 - 5.19* - 5.20* - 5.21*

\begin{ejercicio}[Libro 5.4]
    % //TODO: copiar ejercicio de Jesus

    Sea $(X_i)_{i \ge 1}$ una sucesión v.a. de Bernoulli con $0 < p < 1$. Demuestra que para todo $p < a < 1$ se cumple que:
    \begin{equation}
        \mathbb{P}\left( \frac{1}{n} \sum_{i=1}^{n} X_i \ge a \right) \le e^{-n h(a;p)},
    \end{equation}
    donde
    \begin{equation}
        h(a;p) = a \log \frac{a}{p} + (1 - a) \log \frac{1 - a}{1 - p}.
    \end{equation}
    \textbf{Pista:} Demuestra primero que para todo $s \ge 0$
    \begin{equation}
        \mathbb{P}\left( \frac{1}{n} \sum_{i=1}^{n} X_i \ge a \right) \le e^{-nas} \left( \mathbb{E}[e^{sX_1}] \right)^n.
    \end{equation}
\end{ejercicio}

\begin{ejercicio}[Libro 5.10]
    % //TODO: copiar ejercicio de Jesus
\end{ejercicio}

\section{Sección 8}
% //TODO: ejercicios hacer 8.1 - 8.3 - 8.5 - 8.6 - 8.11

\begin{ejercicio}[Libro 8.1]
    % //TODO: copiar ejecicio de Jesus y Cristian
\end{ejercicio}

\begin{ejercicio}[Libro 8.6]
    % //TODO: copiar ejercicio de Jesus y Cristian
\end{ejercicio}

\section{Sección 9}
% //TODO: ejercicios hacer 9.1 - 9.2 - 9.3 - 9.4 - 9.7 - 9.17

\begin{ejercicio}[Libro 9.2]
    % //TODO: copiar ejercicio de Jesus, Cristian y Ilaria
\end{ejercicio}

\begin{ejercicio}[Libro 9.17]
    % //TODO: copiar ejercicio de Jesus, Cristian y Ilaria
\end{ejercicio}


\section{PDF extra}
% //TODO: ejercicios hacer 1 - 2 - 3 - 9 - 10 - 13
